\section{Whole transcendental lattices under factorization}\label{gw:section}

We construct an operator algebra which distinguishes the whole transcendental
contribution of each surface center.  Its role is to make the classical
integral blowup decomposition compatible with cancellation in an arbitrary
weak factorization.
Quantum blowup decompositions also appear in Iritani's decomposition of
quantum $D$-modules and quantum cohomology $F$-manifolds
\cite[Theorem~1.1 and Corollary~1.2]{Iritani25}.
That decomposition uses a formal Novikov extension and a change of
variables. Here we study the rational operator algebra in the classical
blowup splitting and its saturated integral transcendental summands.

\subsection{Hodge structures and the operator algebra}

For a smooth projective variety $Y$, write
\[
 H^{2i}(Y,\Q)_{\mathrm{Hdg}}
 =H^{2i}(Y,\Q)\cap H^{i,i}(Y).
\]
If $M$ is a rationally connected smooth projective fourfold, set
\[
 L_M=H^4(M,\Z)/\mathrm{tors},\qquad
 T_M=L_M\cap H^4(M,\Q)_{\mathrm{Hdg}}^{\perp}.
\]
The orthogonal is taken for the cup form $q_M$.  For a smooth projective
surface $Z$, use the analogous definition
\[
 T_Z=(H^2(Z,\Z)/\mathrm{tors})\cap
              H^2(Z,\Q)_{\mathrm{Hdg}}^{\perp},
\]
with surface cup form $q_Z$.  We always divide out torsion before discussing
integral lattices.  When comparing with degree four, $T_Z(-1)$ means the Tate
shift to weight four, equipped with the form $-q_Z$.

\begin{lemma}\label{gw:hodge}
The forms on $T_M$ and $T_Z$ are nondegenerate.  The rational Hodge structure
$T_{M,\Q}$ is generated by its $(3,1)$ and $(1,3)$ subspaces, and
$T_{Z,\Q}$ is generated by its $(2,0)$ and $(0,2)$ subspaces.  The only
off-diagonal types in the even cohomology of $M$ are $(3,1)$ and $(1,3)$ in
degree four.  If $i:Z\hookrightarrow M$ is a surface embedding, then
\[
 i^*T_{M,\Q}=0,\qquad i_*T_{Z,\Q}=0.
\]
\end{lemma}
\begin{proof}
Rational connectedness gives $H^{p,0}(M)=0$ for $p>0$
\cite[Theorem~30]{AraujoKollar02}.
We use the Hodge decomposition, functoriality and Lefschetz polarizations
in \cite[Sections~2.1--2.2]{VoisinHodge2016}; the Hodge-complement
statement is Lemma~1 there.
Thus $H^2(M,\Q)$ is of type $(1,1)$; hard Lefschetz gives the analogous
statement in degree six.  The remaining assertion about types follows from
Hodge symmetry.  In degree four all nonprimitive Lefschetz summands are
rational Hodge classes.  Hodge--Riemann therefore makes their orthogonal
primitive transcendental part nondegenerate.  The same argument in degree
two applies to a surface, using an ample class to split off its
nonprimitive part.  Alternatively, the nondegeneracy on the Hodge-class
space follows on each primitive Lefschetz summand from definiteness.

By semisimplicity of polarizable rational Hodge structures, the substructure
generated by the stated off-diagonal types has a Hodge complement in the
transcendental structure.  That complement would have only diagonal type,
and hence would consist of rational Hodge classes.  It is zero by the
definition and nondegeneracy of the transcendental part.  Finally, $i^*$
takes the off-diagonal part of $H^4(M)$ into $H^4(Z)$, which has only type
$(2,2)$; and $i_*$ takes the off-diagonal part of $H^2(Z)$ into $H^6(M)$,
which has only type $(3,3)$.  Both maps vanish on the generating types and
therefore on the whole corresponding rational Hodge structure.
\end{proof}

All Gromov--Witten invariants below are connected, of genus zero, and use
cohomology with rational coefficients.  The \emph{Hodge difference} of a
class of type $(p,q)$ is $p-q$.  All degrees are collected by numerical
curve class: an invariant denotes the sum over homology classes with the specified numerical class.  This
sum is finite in every bounded ample degree, and
all degree comparisons below use an integral ample class.  Other than the
two specified variable insertions, every insertion is a rational class of
diagonal Hodge type.  Descendant powers at all ordinary markings, including
the two variable markings, are arbitrary nonnegative integers.
These two markings carry distinct formal labels in every weighted list
and degeneration graph. The labels remain distinct even when the two
vectors agree or one is zero. Automorphism and gluing coefficients thus
depend on the fixed discrete data, independently of the variable vectors.

\begin{definition}\label{gw:algebra}
For every invariant
\[
 B(u,v)=
 \left\langle\tau_a(u),\tau_b(v),
       \tau_{a_1}(\gamma_1),\ldots,\tau_{a_n}(\gamma_n)
 \right\rangle^M_{0,\beta},\qquad u,v\in T_{M,\Q},
\]
let $A_B\in\operatorname{End}_{\Q}(T_{M,\Q})$ be determined by
$B(u,v)=q_M(A_Bu,v)$.  Let $G_M$ be the unital rational algebra generated by
all such operators.  On the zero space we take the zero algebra.
\end{definition}

The virtual classes, evaluation maps, and cotangent classes defining these
invariants respect Hodge structures.  Thus $A_B$ is a Hodge endomorphism.
Interchanging the two variable markings supplies its adjoint.  In particular,
every expression in $G_M$ is a rational Hodge endomorphism.  By
Lemma~\ref{gw:hodge}, equality of two such expressions can be checked on
$(3,1)$, by pairing their images with $(1,3)$.  We will use the analogous
test on surfaces, after the indicated Tate shift.
Accordingly, when analyzing a tensor on copies of $T_M$ or $T_Z$, we first
test its two inputs on opposite off-diagonal types and then use this
Hodge-generation statement to recover the full rational tensor.

\begin{theorem}[Operator and integral blowup decomposition]\label{gw:blowup}
Let $\pi:Y\to M$ be the blowup of a connected smooth center $Z$ of
codimension $r\geq2$ in a rationally connected smooth projective fourfold.
If $Z$ is a surface, the classical embeddings identify
\begin{equation}\label{gw:split-equation}
(T_Y,q_Y)=(T_M,q_M)\perp(T_Z(-1),-q_Z)
 \quad\text{over }\Z,
 \qquad
 G_Y=G_M\oplus\Q\,\operatorname{id}_{T_Z}.
\end{equation}
The second factor is omitted when $T_Z=0$.  For a point or curve center,
$T_Y=T_M$ integrally and $G_Y=G_M$.
\end{theorem}

Here and subsequently a product decomposition of operator algebras means
that the two factors act independently on the displayed summands, with
zero mixed blocks.  In particular, the assertion is stronger than the
reconstruction of individual numerical invariants.

Retain the blowup $\pi:Y\to M$ of Theorem~\ref{gw:blowup} throughout
its proof. Write $i:Z\hookrightarrow M$ and $j:D\hookrightarrow Y$
for the center and exceptional-divisor embeddings, where
$D=\mathbb P_Z(N_{Z/M})$. Let $p:D\to Z$ and
$\xi=c_1(\mathcal O_D(1))$. Put $W=T_{Z,\Q}$ if $Z$ is a surface,
and $W=0$ otherwise.

We first identify the integral summands in the theorem. The classical
blowup formula has explicit Hodge-compatible integral maps
\cite[Section~2.1, Lemmas~9--10]{Schreieder14}. In the surface case it reads
\[
 H^4(Y,\Z)/\mathrm{tors}
 =\pi^*(H^4(M,\Z)/\mathrm{tors})
       \oplus j_*p^*(H^2(Z,\Z)/\mathrm{tors}).
\]
The summands are orthogonal, since $p_*1=0$, and on the exceptional summand
$j^*j_*=-\xi$ and $p_*\xi=1$ give $-q_Z$. After rationalization this is a
direct sum of Hodge structures, so its Hodge-class subspace splits as well.
Taking the orthogonal and intersecting with the displayed integral direct
sum gives the integral transcendental decomposition. For a curve or
point center the additional degree-four cohomology is diagonal and
contributes no transcendental summand.

The remaining assertion concerns the operator algebra on this fixed
splitting. We first prove
$G_Y\subseteq G_M\oplus\Q\operatorname{id}_W$: there are no mixed
operators, the exceptional block is scalar, and the old block lies in
$G_M$. An exceptional-line invariant then supplies the projector onto
$W$, separating that factor from its complement. Finally we recover
every operator in $G_M$ on the old summand. The two inclusions require
different relative reconstructions: the first uses ordinary insertions
pulled back from $M$, whereas recovery allows arbitrary ordinary
insertions on $Y$. When $W=0$, only the two inclusions are needed.
The next three subsections carry out these steps; the final subsection
uses the product decomposition to cancel whole integral lattices along
a weak factorization.

\subsection{Classical degeneration and local path tensors}

Both reconstruction arguments express a two-variable tensor as contractions
of connected relative tensors. The genus-zero tree reduces these
contractions to the path joining the two variable markings. Along that
path, local projective-bundle vertices supply scalar pairings between
copies of $W$. We first specify the connected-factor and descendant
conventions needed for these statements.

We use expanded relative stable maps in the sense of Jun Li
\cite{Li01,Li02}.  An ordinary marking carries an evaluation class on the
target; a relative marking carries a positive contact order and an
evaluation class on the boundary divisor.  In a degeneration graph the
relative markings are the joining roots.  A descendant at an ordinary marking is the first Chern
class of the cotangent line of the actual source curve there.  These are the
same cotangent lines in every relative theory used below.  Gluing relative
markings to nodes changes no neighborhood of an ordinary marking; hence
these cotangent lines restrict to the corresponding lines on the relative
factors.  On a smooth fibre they are the usual absolute descendant lines.
We never require a descendant at a joining relative marking.

We require the degeneration formula with a product of \emph{connected}
relative factors.  A disconnected source can share one expanded target,
so this product requires a virtual-class theorem.  Here is the precise
chain of results we use.  Apply the ordinary-descendant degeneration formula
of Abramovich--Fantechi \cite[Theorem~0.1 and Section~5]{AF16}.
For fixed labelled connected-component data
$\Xi=\coprod_\nu\Gamma_\nu$, write $\mathcal K_{\Gamma_\nu}(X,D)$ for
the relative-map space with that connected degree, genus, and marking data,
and $\mathcal K_\Xi(X,D)$ for the space with all those components sharing
an expanded target.  Their contraction morphism
\[
 \epsilon:\mathcal K_\Xi(X,D)\longrightarrow
              \prod_\nu\mathcal K_{\Gamma_\nu}(X,D)
\]
separates and stabilizes the components, and satisfies
\[
 \epsilon_*[\mathcal K_\Xi(X,D)]^{\mathrm{vir}}
      =\boxtimes_\nu[\mathcal K_{\Gamma_\nu}(X,D)]^{\mathrm{vir}}
\]
by \cite[Proposition~4.14]{AF16}.  Corollary~4.15 there gives the
ordinary-descendant product formula, retaining all relative evaluation
weights.  This corollary supplies descendant compatibility under the
separating and stabilizing map $\epsilon$.
Its hypothesis that each connected component has a marking
holds here: each vertex of a nontrivial connected degeneration graph has
a joining root.  A one-vertex term needs no product decomposition.
Next apply \cite[Theorem~1.1.2]{AMW14} to each connected factor to identify
it with the Jun Li relative invariant.

These steps preserve our cotangent convention.  The ordinary lines in
\cite[Definition~4.1 and Section~4.3]{AF16} agree with those on the actual
coarse source near ordinary markings.  The comparison morphism from
Abramovich--Fantechi to Li contracts no source components
\cite[Section~4.1]{AMW14}; its virtual pushforward and the projection
formula therefore identify the actual-source descendants as well as the
evaluation classes.  We need no descendant at a joining root and no
comparison for disconnected Li moduli.  This also explains why we use the
product theorem for maps to a fixed relative target; no product rule for
rubber targets is asserted.

The resulting formula expresses an absolute invariant as a finite sum
over matching connected relative data.  Each edge contracts a pair of
relative weights with the diagonal of the double divisor.  The coefficient
is the product of contact orders divided by the automorphism and labeling
factors.  A connected genus-zero source gives a tree of connected vertices.
We retain the two variable insertions throughout, so the formula is an
identity of bilinear tensors.

The two normal-cone degenerations used in the reconstruction have
auxiliary components
\[
 P=\mathbb P_Z(N_{Z/M}\oplus\mathcal O_Z),\qquad
 P_D=\mathbb P_D(N_{D/Y}\oplus\mathcal O_D).
\]
Degenerating $M$ along $Z$ gives the pair $(Y,D)$ joined to $P$ along its
infinity divisor $D$; its zero section is $Z$.
Degenerating $Y$ along $D$ gives $(Y,D)$ joined to $P_D$ along its infinity
section, another copy of $D$; its zero section is also a copy of $D$.
Both auxiliary components, and their expansions, project to $Z$.
All classes of Hodge difference $\pm2$ on $D$, on
$P$, and on the projective bundles
appearing when one degenerates along $D$, come from copies of $W$ multiplied
by tautological classes.  If $W=0$ there are no such classes.

Multiplication of one of these copies of $W$ by a fixed diagonal class
preserves the sum of the copies, and is scalar between its degree-compatible
copies.  Indeed, the bundle formula reduces the assertion to base
multiplication on $Z$; a positive-degree diagonal base class annihilates
$W$, as is seen on $(2,0)$ and $(0,2)$ and then by
Lemma~\ref{gw:hodge}.  A product of two $W$ variables is $q_Z(u,v)$ times
the point class, with any tautological factors retained.  We will use these
facts when tail outputs combine cohomology labels.

\begin{lemma}[Local tensors]\label{gw:local}
On any of the projective-bundle pieces just described, a connected
genus-zero invariant, either absolute or relative to the indicated infinity
boundary, with two variable insertions in the corresponding
copies of $W$ and all other insertions rational and diagonal is a
rational scalar multiple of $q_Z$ on the corresponding copies of $W$.
It vanishes if the projected curve class on $Z$ is nonzero.  These statements
allow ordinary descendants and expanded targets.
\end{lemma}
\begin{proof}
There is nothing to prove if $W=0$.  Otherwise $p_g(Z)>0$, so $Z$ is not
uniruled.  A fixed-degree connected genus-zero stable map with positive
projected degree has projected image a connected union of rational curves.
For each irreducible component of its finite-type moduli space the evaluation
image in $Z$ has dimension at most one: a dominating evaluation image would
give a covering family of rational curves on $Z$.  There are finitely many
such components for the fixed discrete data.  A holomorphic two-form pulls
back to zero through each of their evaluation maps, since their images lie
in curves or points.  The bilinear tensor, with all other data fixed, is
rational and respects Hodge structures.  Its associated Hodge endomorphism
therefore annihilates the $(2,0)$ generating subspace, and its conjugate,
and hence all of $W$ by Lemma~\ref{gw:hodge}.  This proves vanishing as a
tensor, independently of the values of the two variables.  The argument
also applies after projecting an expanded target to $Z$.

If the projected degree is zero, all evaluations in $Z$ agree.  Factoring out
base classes by the projection formula, the two variable classes occur as
$u\smile v$.  They already have top degree on the surface.  Every additional
positive-degree base factor kills their product; the remaining fibre
integral is a rational number multiplying $\int_Zuv$.  This reasoning is
independent of $u,v$ and allows the specified descendants.
\end{proof}

\begin{lemma}[A tree with two variables]\label{gw:path}
In either normal-cone degeneration above, consider a genus-zero tree with
exactly two variable insertions of
Hodge differences $+2$ and $-2$.  Every branch disjoint from the path joining
them supplies a rational diagonal class at its attachment.  Along that path
the variable spaces are copies of $T_{M,\Q}$ or $W$, with the appropriate
Tate shifts.  Orient the path from the first labelled variable marking to
the second, and order the two slots at each path vertex in that direction.
Suppose that, with this slot order, each vertex tensor belongs to the
applicable row of the following table.  Then contraction of the path
belongs to the row determined by its two remaining ends:
\[
 \begin{array}{c|c}
 \text{two end spaces}&\text{allowed tensors}\\ \hline
 T_{M,\Q},T_{M,\Q}
   & (u,v)\longmapsto q_M(Au,v),\quad A\in G_M,\\
 W,W&\Q q_Z,\\
 T_{M,\Q},W\text{ or }W,T_{M,\Q}&0.
 \end{array}
\]
One may replace $G_M$ by any specified rational unital operator algebra,
provided that each path vertex still satisfies its applicable row in this
oriented slot order.
In the reconstruction arguments below, this vertex hypothesis is supplied
by the local-tensor lemma or by an earlier induction case.
\end{lemma}
\begin{proof}
Cut an edge outside the path.  The detached connected subtree has only
rational diagonal input classes.  Its single output is rational and of
Hodge difference zero, because its virtual class and all structure maps
respect Hodge difference.  At the retained vertex the attachment remains
its existing relative root,
with the same contact and the resulting diagonal evaluation weight; it is
not an additional ordinary marking.  Repeating this operation removes all
side branches from the tensor contraction.  Conservation of Hodge difference
at the remaining vertices forces difference two along the path.  Lemma~\ref{gw:hodge} and the
projective bundle formula identify the indicated variable spaces.
These path classes and the diagonal branch outputs have even cohomological
degree, so their permutations introduce no odd-degree signs.
For the first row, contraction of consecutive tensors written in path order
as $q_M(Au,x)$ and $q_M(By,v)$ pairs $x$ and $y$ with the inverse cup form
and gives $q_M(BAu,v)$.  Thus contraction composes operators without taking
adjoints.  In the second row it composes scalar maps between copies of $W$.
A path changing between the two sorts
has a mixed vertex, whose tensor vanishes by the stated vertex hypothesis.
The diagonal of a projective bundle gives
the same identities between its copies of $W$, with the relevant nonzero
pairing constants.  The conclusion follows from closure under composition
and rational linear combination.  If the two variables meet on one local
piece, Lemma~\ref{gw:local} supplies the second row directly.
\end{proof}

\subsection{Reconstruction with a center of codimension at least two}

We first prove a relative assertion for $(Y,D)$.  Ordinary insertions are
pullbacks from $M$.  Relative weights are written
\begin{equation}\label{gw:relative-weights}
\mu=\bigl((k_j,\xi^{b_j}p^*\delta_j)\bigr)_{j=1}^{L},\qquad k_j\geq1,
       \quad0\leq b_j\leq r-1.
\end{equation}
The two variables may be ordinary $T_M$ insertions or relative $W$ insertions.
All other classes are rational and diagonal.  A relative invariant with
these data is interpreted as a bilinear tensor on its two variable spaces.
Write $B_{\mathrm{rel}}(u,v)$ for such a tensor in numerical curve class
$\beta$ on $Y$, with relative list $\mu$.  Admissibility requires
$D\cdot\beta=\sum_{j=1}^{L}k_j$.
The \emph{original ordinary markings} are the ordinary markings of this
relative datum, before translating any relative conditions into additional
absolute markings.  This distinction fixes the ordinary-mark count used
in the induction.

The translation of relative conditions into absolute descendants follows
Hu--Li--Ruan \cite[Section~5.2, equation~(17)]{HLR08}, extending the divisor
correspondence of Maulik--Pandharipande
\cite[Theorem~2 and Section~2.4]{MP06}.
The induction below retains the two variable insertions and all original
ordinary descendants, in order to prove the required tensor restrictions.

\begin{lemma}[Restricted relative reconstruction]\label{gw:relative}
Every such tensor satisfies the three rows of Lemma~\ref{gw:path}.
\end{lemma}
\begin{proof}
Replace a relative weight $(k,\xi^bp^*\delta)$ by the ordinary insertion
\begin{equation}\label{gw:translation}
\tau_{r(k-1)+b}(i_*\delta)
\end{equation}
to form the absolute tensor $B_{\mathrm{abs}}(u,v)$ on $M$ in numerical
class $\pi_*\beta$, retaining all original ordinary insertions
and their descendant powers.  We apply degeneration to the normal cone
using the following specific lifts.  In
\[
 \mathcal W=\operatorname{Bl}_{Z\times\{0\}}
                     (M\times\mathbb A^1)
\]
the strict transform $\mathcal Z$ of $Z\times\mathbb A^1$ is still
$Z\times\mathbb A^1$, because the blowup ideal restricts to the Cartier ideal
$(t)$.  Lift $i_*\delta$ by $(i_{\mathcal Z})_*(\delta\times1)$.
On the special fibre its restriction to $Y$ is zero and its restriction to
$P=\mathbb P_Z(N_{Z/M}\oplus\mathcal O_Z)$ is the zero-section Gysin class.
These are transverse fibre restrictions.  Original ordinary classes use
pullbacks from $M$.

This lift is valid even when $i_*\delta=0$.  A class on $\mathcal W$ may
restrict to zero on the smooth fibre and nontrivially to its special
components.  For a relative variable in $W$, Lemma~\ref{gw:hodge} says that
the associated absolute insertion is zero, so the absolute bilinear tensor
is zero.  For two ordinary $T_M$ variables the associated absolute tensor
is a generator allowed in the first row.  Thus the absolute starting tensor
has the required restriction in all three cases.

We prove the relative restriction by simultaneous strong induction, keeping
the two variables throughout.  Fix an integral ample divisor $H$ on $M$.
The outer order is first $H\cdot\pi_*\beta$, then the number of original
ordinary markings.  Within fixed outer data, use, in succession, the length
of the relative list, the sum of its contacts, and the sum of its
hyperplane exponents, with smaller values first.  At every smaller outer
index the assertion is assumed simultaneously for all numerical curve
classes at that index, relative lists, evaluation labels, placements of
the two formal variable slots, and ordinary descendant exponents.
At equal outer indices the inner induction is likewise simultaneous in
the curve classes, labels, variable placements, and descendant exponents.
Cutting a side branch only replaces the weight
at an existing relative root, so it does not increase the number of original
ordinary markings.

A connected relative vertex on $Y$ of zero projected degree cannot have
joining markings.  Indeed every nonzero effective curve contracted by
$\pi$ has class a positive multiple of an exceptional line, and has negative
$D$-degree; relative admissibility would require its $D$-degree to equal a
sum of positive contacts.  A zero-class vertex likewise has no contacts.
Thus zero projected degree contributes only a constant invariant with no
joining marks.  Its two middle-degree ordinary variable classes already
have total top degree; the invariant is a scalar cup pairing, or zero.
This supplies the initial case.

For a noninitial term of the degeneration formula, cut away the branches
with no variables, as in Lemma~\ref{gw:path}.
Every $Y$-vertex in a nontrivial connected tree has a joining edge, so the
zero-degree argument above makes its projected ample degree positive.
The projected classes of all components are effective and add to the
original class.  Hence, if there is more than one $Y$-vertex, every one
has smaller projected degree and its path tensor is controlled by induction.
A unique $Y$-vertex of smaller projected degree is controlled in the same way.
If a $Y$-vertex has the full projected degree, it is the unique $Y$-vertex;
its tensor is still an earlier case if it retains fewer original ordinary
markings.  Local path vertices are controlled by Lemma~\ref{gw:local}.
In these cases every path vertex therefore satisfies its row, so the
conditional contraction rule of Lemma~\ref{gw:path} controls the term.
A path containing no $Y$-vertex is controlled entirely by the local lemma.

It remains to consider a unique $Y$-vertex of the full projected class
carrying all original ordinary markings.
Every $P$-vertex is then a fibre tail with one joining marking, since the
graph is a tree.  Let such a tail have joining contact $k'$ and receive the
translated markings indexed by a set $J$.  All its evaluations in $Z$
agree.  Its pushforward to the joining divisor therefore has the form
\[
 \left(\prod_{j\in J}p^*\delta_j\right)\eta,
\]
where $\eta$ is a rational diagonal class independent of the base labels.
The projection formula gives this identity also when a label is one of the
two variables.  The complex codimension of $\eta$ is
\begin{equation}\label{gw:tail-codim}
c=r\left(\sum_{j\in J}k_j-k'\right)
       +\sum_{j\in J}b_j-|J|+1.
\end{equation}
Indeed the relative fibre virtual dimension is
$r-2+rk'+|J|$; each translated zero-section insertion has fibre codimension
$r+r(k_j-1)+b_j=rk_j+b_j$, and pushing to the relative fibre divisor of
dimension $r-1$ gives the displayed difference.

If one tail contains both variable markings, both are relative $W$
variables, since all original ordinary markings remain on the main vertex.
Their product is $q_Z(u,v)$ times the point class on $Z$, so the tail output
is $q_Z(u,v)$ times a fixed class.  The remaining vertices then have only
fixed insertions, so the graph is a rational multiple of $q_Z$.  It already
lies in the allowed $W,W$ row, without using induction on its main tensor.
For every other nonzero term, expand the tail outputs in the projective
bundle basis.  A tail with one $W$ variable multiplies it only by fixed
diagonal classes, which act by scalars between the allowed copies or
annihilate it.  Thus each remaining summand retains two separate variable
slots in the prescribed spaces, and the following induction applies to it.

A negative $c$ gives zero.  An empty tail has $c=1-rk'<0$, so every tail
receives a translated marking.  The new relative list therefore has no
more entries than the old one.  A strict decrease is an earlier induction
case.  If the lengths agree, each tail has exactly one translated marking.
Then $c\geq0$ and $b_j<r$ imply $k'\leq k_j$.  A strict decrease of the
total contact is again earlier.  At equal contacts, expand $\eta$ in the
projective bundle basis.  Its fibre exponent is at most $b_j$, so either
the total hyperplane exponent decreases or all the original indices agree.
A decrease of a fibre exponent only multiplies its base label by a rational
diagonal class, and preserves the stated variable spaces.  Thus all smaller
terms with separate variable slots retain the required tensor
restriction.

For one translated mark and equal indices, the coefficient of the highest
fibre power is obtained by restricting to a fibre.  The relative invariant
computed in \cite[Theorem~7.1]{HLR08} is
\begin{equation}\label{gw:fibre-number}
\left\langle\tau_{rk-1-j}(\mathrm{pt})\mid H^j\right\rangle^%
 { (\mathbb P^r,\mathbb P^{r-1})}_{0,k[\mathrm{line}]}
 =\frac{1}{k^{r-j}((k-1)!)^r}>0,
 \qquad 0\leq j\leq r-1.
\end{equation}
The proof of that theorem identifies the descendant line with the
cotangent line of the source at the ordinary marked point, so its convention
is the one fixed above.  Taking $j=r-1-b$ gives exactly the descendant in
\eqref{gw:translation}.  Gluing multiplies these numbers by positive contact
and labeling factors.  Consequently the remaining term is a nonzero
rational scalar times the original relative tensor.  At fixed projection,
the total contact fixes the numerical class on the blowup, since the
kernel on curve classes is generated by the exceptional line.

We have an identity $B_{\mathrm{abs}}=C B_{\mathrm{rel}}+R$, with
$C\in\Q^*$ independent of the two variable vectors.  The starting tensor
and the remainder belong to the relevant row of Lemma~\ref{gw:path}; hence
so does $B_{\mathrm{rel}}$.  The induction is simultaneous in the three
rows.  It is a tensor identity, not a separate choice of a scalar for each
pair of vectors.  Testing on opposite off-diagonal types and using
Lemma~\ref{gw:hodge} extends it to the full rational transcendental
structures.  This completes the induction.
\end{proof}

\subsection{The exceptional factor and recovery of the old algebra}

We first obtain
\begin{equation}\label{gw:inclusion}
G_Y\subseteq G_M\oplus\Q\operatorname{id}_{W}.
\end{equation}
Degenerate $Y$ along its divisor $D$.  The main component is again $Y$,
and the other component is the projective-line bundle $P_D\to D$ defined above.
By the blowup cohomology formula every ordinary insertion on $Y$ is a sum
of a pullback from $M$ and exceptional terms $j_*a$, where $j:D\hookrightarrow
Y$.  Lift an exceptional term by the strict transform of the divisor family,
so that it is supported on the zero section of the bundle component.
Use pullbacks for the other terms.  Thus the ordinary insertions on the
main relative component are exactly those of Lemma~\ref{gw:relative}.
On the other components the local rule of Lemma~\ref{gw:local} applies,
through their projection to $Z$.  Applying Lemma~\ref{gw:path} to each
degeneration tree proves \eqref{gw:inclusion}, including zero mixed blocks.

The inclusion controls every operator but does not yet separate the two
factors or recover all operators on the old summand. The exceptional
line gives the required separation.
If $W\neq0$, let $e$ be the exceptional line class.  Every stable map of
class $e$ is contained in one fibre of $D\to Z$.  A fibre has normal bundle
$\mathcal O(-1)\oplus\mathcal O^{\oplus2}$ in $Y$, so these maps are
unobstructed.  The two-mark degree-one evaluation map on the fibre has
degree one onto $\mathbb P^1\times\mathbb P^1$.  Restricting
$j_*p^*u$ to $D$ gives $-\xi p^*u$, whence
\begin{equation}\label{gw:exceptional-line}
\left\langle j_*p^*u,j_*p^*v\right\rangle^Y_{0,e}
       =\int_Zuv.
\end{equation}
An insertion from $T_M$ restricts to zero on $D$ by
Lemma~\ref{gw:hodge}.  Relative to the negative cup form of the exceptional
summand, \eqref{gw:exceptional-line} is minus its projector.  Thus the
projector onto $W$ belongs to $G_Y$.  Its complement belongs to $G_Y$ too.

It remains to recover $G_M$.  For this purpose we need the divisor version
of Lemma~\ref{gw:relative}, now allowing all ordinary classes on $Y$.
Write $\operatorname{pr}_M G_Y$ for its algebra of restrictions to $T_M$;
this is defined by \eqref{gw:inclusion}.

\begin{lemma}[Divisor reconstruction with separated variable spaces]
\label{gw:divisor}
For $(Y,D)$, allow arbitrary ordinary diagonal insertions and arbitrary
ordinary descendants.  Allow the two variables either as ordinary classes
in $T_M$ or $W\subset H^4(Y)$, or as boundary classes in a tautological
copy of $W$ in $H^*(D)$.  Then relative tensors with two $T_M$ ends belong
to $\operatorname{pr}_M G_Y$, tensors with two $W$ ends are scalar cup
pairings, and mixed tensors vanish.
\end{lemma}
\begin{proof}
Fix a relative tensor $B_{\mathrm{rel}}(u,v)$ in numerical curve class
$\beta$ on $Y$, with relative list
$\mu=((k_j,\delta_j))_{j=1}^{L}$.  Thus
$D\cdot\beta=\sum_{j=1}^{L}k_j$.
Use degeneration of $Y$ along $D$, with the same source cotangent convention,
and translate a relative weight $(k,\delta)$ into
$\tau_{k-1}(j_*\delta)$ to form an absolute tensor
$B_{\mathrm{abs}}(u,v)$ on $Y$ in the same numerical class $\beta$,
retaining the original ordinary insertions and descendants.
Use the strict divisor-family lift for translated
classes and pullbacks for original ordinary classes.  This is the
classical divisor correspondence of \cite[Section~2.4, Lemma~4]{MP06}; we
give the ordering needed for the restricted tensors.

The ordinary $T_M$ part restricts to zero on $D$.  The rational Hodge
structures generated by boundary classes of Hodge difference $\pm2$ are
copies of $W$: for a surface center they are
$p^*W\subset H^2(D)$ and $\xi p^*W\subset H^4(D)$.
Gysin carries the first into the exceptional $W$ summand of $H^4(Y)$ and
annihilates the second, since $H^6(Y)$ is diagonal.  Restriction of the
ordinary exceptional class is $j^*j_*p^*u=-\xi p^*u$.
It follows from \eqref{gw:inclusion} that every associated absolute tensor
has the claimed restriction.  Local tensors have it by
Lemma~\ref{gw:local}.  This also treats a translated variable with zero
Gysin class; the absolute tensor is then zero.

Fix an integral ample divisor on $Y$ and perform outer strong induction
on its curve degree, then on the number of original ordinary markings.
As before, these are the ordinary markings before translation, and the
induction is simultaneous for all numerical classes, relative data,
labels, variable placements, and descendant exponents at the given index.
Pruning replaces existing relative weights as in Lemma~\ref{gw:path}.

Every main $Y$-vertex in a nontrivial tree has a joining root.  Its curve
class cannot be zero, since a zero class has total contact zero; its ample
degree is therefore positive.  Under collapse of the degeneration to $Y$,
the projected component classes are effective and add to the original
class $\beta$.  If there are several main vertices, all have smaller ample
degree.  A unique main vertex of smaller degree is also an earlier case.
A main vertex with full degree is unique, and is an earlier case
unless it retains all original ordinary markings.  Thus in these cases
induction controls every main path vertex and Lemma~\ref{gw:local}
controls every local one, as required by the conditional path rule.
A path with no main vertex is controlled by the local lemma after pruning.

At equal outer data the unique main vertex has curve class $\beta$:
all other components project to points of $Y$.  They are fibre tails of
$P_D\to D$, since the graph is a tree.  This argument uses the ample degree
on $Y$, without requiring $D$ to be nef.  Hence
$K=D\cdot\beta$, the total contact, is fixed.  If $K<0$ there is no
admissible relative datum.  A zero curve class has $K=0$ and no relative
marks; its constant tensors are scalar cup pairings, zero on mixed
summands.  More generally, for $K=0$ the following inner induction has just
the empty-list case.

Let the new main relative list have length $t$, with tail contacts $k'_a$.
If $J_a$ is the set of
translated markings on tail $a$, the tail calculation
\eqref{gw:tail-codim} now applies with $r=1$: the center of this
normal-cone degeneration is the divisor $D\subset Y$.
Thus
\begin{equation}\label{gw:divisor-order}
c_a=\sum_{j\in J_a}(k_j-1)-k'_a+1,
 \qquad
 \sum_a c_a=t-L.
\end{equation}
Before applying the list order, use the tail-output factorization from the
first reconstruction.  A tail containing both variable slots gives
$q_Z(u,v)$ times fixed classes, so its whole graph contribution is already
an allowed scalar pairing.  In every other nonzero summand the variable
slots remain separate: fixed diagonal classes act by scalars between the
copies of $W$.  The following inner induction applies to these summands.

Each nonzero tail has $c_a\geq0$.  Therefore $t\geq L$.
Use decreasing length as the next order; it terminates because
$t\leq K$.  This is an inner induction for each fixed $\beta$, so no
unbounded reversal of an order is being used.

At equal lengths every $c_a$ is zero.  Each tail has
\[
 k'_a-1=\sum_{j\in J_a}(k_j-1),\qquad
 \delta'_a=\text{a scalar times }\prod_{j\in J_a}\delta_j.
\]
Take all fixed labels homogeneous.  Give $(k_j,\delta_j)$ weight
$k_j-1+\operatorname{codim}\delta_j$.  These weights are nonnegative.
Each formal variable slot has a fixed cohomological degree, unchanged
when its vector is specialized to zero or to the other vector.  A boundary
$W$ slot has positive codimension, so it always has positive weight.
Combining two positive-weight labels on a tail decreases their number, so
use that number, with smaller values first, as the final order.
A strict decrease of this count is an earlier case with two separately
labelled variable slots.
If the number of positive-weight labels is unchanged, each tail has at most
one positive-weight label.  All other labels are $(1,1)$, so every nontrivial contact and label
is unchanged.  Empty tails have degree one and identity output; redistribution
of the trivial labels, at the fixed total length, consequently gives the
original weighted partition, with each formal variable slot still
distinguished.

These unchanged-data terms cannot cancel.  A tail carrying its one
positive-weight label has the leading fibre invariant in
\eqref{gw:fibre-number}, now for $r=1$, with its positive gluing contact
factor.  An additional $(1,1)$ label translates to an ordinary point-divisor
insertion on $\mathbb P^1$.  By the divisor equation it multiplies the
leading number by $k'_a$.  The descendant correction contains the square
of the fibre point class and vanishes.  For the source cotangent convention
this is the ordinary relative divisor equation with the divisor supported
in the interior, disjoint from the relative boundary.  A degree-one tail
without translated markings contributes one.  Thus every unchanged-data
coefficient is a product of the same positive leading numbers and positive
contact, divisor, and labeling factors.  All path labels are even, so no
permutation signs occur.  The graph with one leading fibre tail for each
relative entry is among these terms.  Their sum is therefore a nonzero
rational coefficient $C$, depending only on the fixed discrete data and
the formal labels, not on the two variable vectors.  For the empty list
the coefficient is one.  This is the nonvanishing part of the triangular
coefficient calculation in \cite[Section~2.4, Lemma~4]{MP06}; its exact
normalization is not needed here.

As before, the formula is an identity of bilinear tensors.  All earlier
terms are allowed by the simultaneous induction and the path lemma; division
by the nonzero, vector-independent coefficient proves the three stated
restrictions.  Ordinary descendants at retained markings have not changed,
and those at markings moved to a fibre occur only in the smaller ordinary-mark
cases.  Hence no bound or extra hypothesis on original descendants or
ordinary diagonal insertions has been used.
\end{proof}

To finish Theorem~\ref{gw:blowup}, degenerate any absolute generator on $M$
along $Z$.  Its two $T_M$ insertions have zero restriction to the normal
bundle component, so they occur on the main side.  By
Lemma~\ref{gw:divisor}, a main vertex carrying both insertions supplies
a tensor in $\operatorname{pr}_M G_Y$ after its side branches are pruned.
If the two insertions lie on different main vertices, the path leaves
the first through a boundary copy of $W$.  That vertex has one $T_M$ end
and one $W$ end, so its tensor vanishes by the mixed row of the same lemma.
When $W=0$, there is no Hodge-difference-two boundary class through which
such a path could pass.  Local tensors and branch outputs are controlled
by Lemmas~\ref{gw:local} and~\ref{gw:path}.  Thus each term, and hence
this absolute generator, belongs to
$\operatorname{pr}_M G_Y$.  Together with \eqref{gw:inclusion} this gives
$G_M=\operatorname{pr}_M G_Y$.  The exceptional-line projector separates
the other factor when it exists.  This proves the theorem.

\subsection{Integral cancellation and the surface forced by rationality}

\begin{lemma}[Whole blocks]\label{gw:blocks}
Along a smooth weak factorization starting at $\mathbb P^4$, each
transcendental lattice is an integral orthogonal sum of whole shifted
surface-transcendental lattices.  Its operator algebra is the product of
independent rational scalar algebras on those summands.  A blowdown with
nonzero surface-transcendental contribution removes one whole summand,
whose integral polarized Hodge-isometry class is that of the entire
transcendental lattice of its center.  A step with zero transcendental
contribution leaves the decomposition unchanged.
\end{lemma}
\begin{proof}
At $\mathbb P^4$ the transcendental lattice is zero.  A blowup with a
nonzero surface contribution adds exactly one summand by
Theorem~\ref{gw:blowup}; a blowup with no transcendental contribution does
nothing.  A blowdown with zero transcendental contribution also leaves the
lattice and algebra unchanged by the same theorem.  Suppose instead that
the next step is a blowdown with nonzero surface contribution.  The theorem,
read in reverse, exhibits the removed center as a factor $\Q$ of
the operator algebra of the current variety.  Its identity is a primitive
central idempotent.  The primitive central idempotents of the already
identified algebra $\Q^m$ are precisely the projectors onto its existing
whole summands.  Thus the removed rational subspace is one of those
summands, even if some summands are mutually Hodge-isomorphic or a surface
transcendental structure is reducible.

The lattice of each summand is its intersection with the ambient integral
lattice: the historical decomposition and the geometric blowup decomposition
are integral direct sums.  Equality of their rational subspaces therefore
identifies their saturated integral lattices and cup forms.  The residual
summands are unchanged.  This proves the induction.
\end{proof}

The numerical comparison below is the rank-twenty-one counterpart of
the surface-classification argument in \cite[Lemma~3]{Kulikov08}.

\begin{lemma}[Rank twenty-one centers]\label{gw:k3-center}
Let $Z$ be a smooth connected projective surface with
$h^{2,0}(Z)=1$ and $\operatorname{rank}T_Z=21$.  Then $Z$ is birational to a
Picard-number-one K3 surface.  Its primitive ample generator has square
$|\operatorname{disc}T_Z|$.  If this square is $d>2$ and
$\operatorname{End}_{\mathrm{Hdg}}(T_{Z,\Q})=\Q$, that K3 surface has
trivial automorphism group.
\end{lemma}
\begin{proof}
Geometric genus and the transcendental lattice are unchanged by point
blowups. Pass to a minimal model using \cite[Proposition~II.16]{Beauville78}.
In the following numerical calculation, $K$, $b_2$, and $q=h^{1,0}$
refer to that minimal surface.
Since $p_g=1$, it is not ruled and has nonnegative Kodaira dimension
\cite[Propositions~III.20--III.21 and Example~VII.3]{Beauville78}. Its canonical class is nef
\cite[Corollary~VI.18]{Beauville78}, so $K^2\geq0$.
Noether's formula \cite[Sections~I.14 and~III.19]{Beauville78} gives
$c_2=12(1-q+p_g)-K^2=24-12q-K^2$.
On the other hand its topological Euler characteristic is
$c_2=2-4q+b_2$.  Hence
\[
                     b_2=22-8q-K^2.
\]
Projectivity gives Picard rank at least one, while the rank assumption gives
$b_2\geq22$.  Thus $q=0$, $K^2=0$, and the Picard rank is one.
Kodaira dimension two is excluded by $K^2=0$ \cite[Proposition~X.1]{Beauville78}.  In Kodaira dimension one the
elliptic fibration \cite[Proposition~IX.2]{Beauville78} supplies a
nonzero isotropic divisor class, independent
of an ample class, contradicting Picard rank one.  The classification in
Kodaira dimension zero leaves a K3 surface, since $p_g=1$ and $q=0$
\cite[Theorem~VIII.2]{Beauville78}.

The K3 lattice is unimodular, and its N\'eron--Severi lattice is primitive
by the integral Lefschetz $(1,1)$ theorem
\cite[Theorem~2]{VoisinHodge2016}.
For a primitive sublattice of a unimodular lattice, its discriminant group
and that of its orthogonal have equal orders.  The N\'eron--Severi lattice
here is generated by a primitive ample class $L$, so its discriminant is
$L^2$.  This proves the degree assertion.

Any automorphism fixes $L$ and hence acts trivially on this rank-one
N\'eron--Severi lattice.  Its action on $T_{Z,\Q}$ is a rational scalar by
the endomorphism hypothesis, and preservation of the nondegenerate form
makes that scalar $+1$ or $-1$.  The unimodular gluing identifies the
cyclic discriminant group of $T_Z$ with that of $\Z L$, of order $d$;
see \cite[Proposition~1.6.1]{Nikulin80} for the compatible gluing of
primitive orthogonal sublattices of an even unimodular lattice.
The action $-1$ on the first and $+1$ on the second is compatible only if
$2$ annihilates this cyclic group, contrary to $d>2$.  Thus the action is
$+1$ on both lattices, and hence on the whole integral second cohomology.
The faithful cohomological action for K3 automorphisms, a consequence of
global Torelli \cite[Proposition~15.2.1]{Huybrechts16}, makes the
automorphism the identity.
\end{proof}

\begin{proposition}[The contribution forced by rationality]
\label{gw:net-center}
Let $M$ be a rationally connected smooth projective fourfold with
\[
 h^{3,1}(M)=1,\qquad \operatorname{rank}T_M=21,\qquad
 |\operatorname{disc}T_M|=d>2,\qquad
 \operatorname{End}_{\mathrm{Hdg}}(T_{M,\Q})=\Q.
\]
Let $\mathcal S$ consist of the Picard-number-one K3 surfaces $S$ whose
shifted integral transcendental lattice $(T_S(-1),-q_S)$ is Hodge-isometric
to $(T_M,q_M)$, and use this collection in the test $u$ of
the preliminary section.  If $f:\mathbb P^4\dashrightarrow M$ is a
birational map, then
\[
                              c_u(f)=1.
\]
Every member of $\mathcal S$ has degree $d$ and trivial automorphism group.
\end{proposition}
\begin{proof}
Under the rationality assumption all varieties in the smooth factorization
of Theorem~\ref{pre:weak-factorization} are rationally connected.
Lemma~\ref{gw:blocks} tracks integral whole blocks through the factorization.
A nonzero transcendental rational Hodge structure of this kind has a
nonzero off-diagonal part by Lemma~\ref{gw:hodge}.  Since the final
$h^{3,1}$ is one, its transcendental structure is simple: a nontrivial
semisimple decomposition would require at least two nonzero $(3,1)$ parts.
Thus the final lattice is one whole block.  The number of additions minus
removals of blocks in its integral polarized Hodge-isometry class is one.

Every center in this class has geometric genus one, transcendental rank
$21$, determinant of absolute value $d$, and rational Hodge endomorphism
ring $\Q$.  Lemma~\ref{gw:k3-center} makes it birational to a member of
$\mathcal S$, and proves the degree and automorphism assertions.  Conversely,
a smooth surface center counted by $u$ is birational to a member of
$\mathcal S$: in dimension two an MRC base which is a K3 has the same
function field as the surface.  Its entire transcendental lattice therefore
has exactly the required class.  Point and curve centers are not counted.

The exceptional divisor of a smooth blowup is a projective bundle over its
center and has the same MRC base, by Lemma~\ref{pre:mrc-fibration}.
The sign convention and additivity in Lemma~\ref{pre:cocycle}
therefore identify $c_u(f)$ with the above number of block additions minus
removals.  This number is one.
\end{proof}
